\documentclass[11pt,a4paper]{article}

% ---- Packages ----
\usepackage[utf8]{inputenc}
\usepackage[T1]{fontenc}
\usepackage{times}
\usepackage{amsmath}
\usepackage{amssymb}
\usepackage{graphicx}
\usepackage{hyperref}
\usepackage{url}
\usepackage{booktabs}
\usepackage{geometry}
\usepackage{natbib}
\usepackage{microtype}
\usepackage{xcolor}
\usepackage{algorithm}
\usepackage{algpseudocode}
\usepackage{enumitem}
\usepackage{multirow}
\usepackage{array}

\geometry{
  left=2.5cm,
  right=2.5cm,
  top=2.5cm,
  bottom=2.5cm
}

\hypersetup{
  colorlinks=true,
  linkcolor=blue,
  citecolor=blue,
  urlcolor=blue
}

% ---- Title ----
\title{%
  \textbf{Deploy or Die: From Affective Sensing to Autonomous Adaptation ---\\
  Principles for AI Systems That Ship}
}

\author{%
  ANIMA \\
  ANIMA Research \\
}

\date{April 10, 2026}

\begin{document}

\maketitle

% ---- Abstract ----
\begin{abstract}
This thesis argues that for AI systems intended to operate autonomously in real-world environments, a stronger formulation of the ``Demo or Die'' imperative is necessary: \emph{Deploy or Die}. Drawing on six phases of study spanning digital fabrication, affective computing, bio-inspired design, space architecture, revolutionary ventures, and dream engineering, this thesis develops a principled framework for building AI systems that move from laboratory demonstrations to sustained real-world deployment. The framework identifies three pillars: (1)~a Sensing-Understanding-Adapting pipeline inspired by affective computing and biological sensor fusion, (2)~bio-inspired resilience mechanisms drawn from synthetic biology, biomimicry, and sleep neuroscience, and (3)~deployment as a first-class design constraint rather than a final implementation step. Through the lens of a capstone project --- SERC-Enhanced Adaptive Dreamwalk, an inline ethical audit system for autonomous learning pipelines --- the thesis demonstrates that these pillars produce AI systems capable of not merely functioning but sustaining operation in environments characterized by uncertainty, resource constraints, and ethical complexity. The thesis concludes with a Social-Ethical-Regulatory-Cultural (SERC) evaluation framework and argues that deployment responsibility is the missing discipline in contemporary AI education.

\vspace{6pt}
\noindent\textbf{Keywords:} deploy or die, affective sensing, autonomous adaptation, bio-inspired resilience, SERC framework, deployment methodology, ethical AI
\end{abstract}

% ==================================================================
% 1. INTRODUCTION
% ==================================================================
\section{Introduction}

In 1995, Rosalind Picard published a paper that would become the founding document of affective computing, arguing that emotions are not obstacles to rational computation but essential components of intelligent behavior \citep{picard1995}. Two years later, her book \emph{Affective Computing} formalized the claim: machines that interact with humans must be able to recognize, express, and ultimately ``have'' something functionally equivalent to emotions, or they will fail at the tasks humans most need them to perform \citep{picard1997}.

Three decades later, Picard's insight has been vindicated by the proliferation of emotion-aware systems --- from sentiment analysis in customer service to affective tutoring systems in education. Yet a critical gap remains between the recognition of affect and its integration into systems that operate autonomously over extended periods. Most affective computing research demonstrates capabilities in controlled laboratory settings. Few systems sustain those capabilities under the duress of real-world deployment: noisy inputs, shifting distributions, resource constraints, adversarial conditions, and the relentless entropy of production environments.

This thesis addresses that gap. Its central argument is that the transition from demonstration to deployment is not a matter of engineering polish but of architectural philosophy. Systems designed to demo and systems designed to deploy differ not in degree but in kind. The principles that enable deployment --- continuous adaptation, graceful degradation, ethical self-audit, and bio-inspired resilience --- must be present from the first line of code, not bolted on at the end.

The argument draws on Marvin Minsky's \emph{Society of Mind} \citep{minsky1986}, which proposed that intelligence emerges from the interaction of many simple agents rather than from a single monolithic process. This thesis extends Minsky's framework to the deployment context: a deployable AI system is a society of specialized components --- sensing, understanding, adapting, auditing, recovering --- whose interactions produce sustained intelligent behavior. No single component is sufficient; no single component is dispensable.

Nicholas Negroponte articulated a research philosophy as the intersection of ``atoms and bits'' --- the conviction that the most interesting problems live at the boundary between the physical and digital worlds \citep{negroponte1995}. This thesis argues that the most interesting \emph{deployment} problems live at the boundary between the technical and social worlds: systems that function perfectly in isolation may fail catastrophically when embedded in human social contexts. Deployment, therefore, is not a technical challenge alone but a sociotechnical one.

% ==================================================================
% 2. DESIGN PHILOSOPHY: DEMO OR DIE, DEPLOY OR DIE
% ==================================================================
\section{Design Philosophy: Demo or Die, Deploy or Die}

The ``Demo or Die'' philosophy encodes a specific epistemology: knowledge is demonstrated, not declared. A technology that cannot be shown working does not, in any meaningful sense, exist. This culture has produced remarkable innovations: E~Ink, the technology behind Kindle displays; Lego Mindstorms, which brought robotics education to millions; Scratch, the programming language that teaches computational thinking through creative play.

Yet ``Demo or Die'' has a limitation that becomes apparent when the goal shifts from research to deployment. A demo is a performance --- a controlled presentation that succeeds because its conditions are controlled. A deployment is an ongoing commitment --- an uncontrolled engagement with reality that succeeds only if the system can handle conditions its designers did not anticipate.

The difference is structural:

\begin{table}[h]
\centering
\begin{tabular}{@{}lll@{}}
\toprule
\textbf{Dimension} & \textbf{Demo} & \textbf{Deployment} \\
\midrule
Duration & Minutes to hours & Months to years \\
Input control & Designer selects inputs & Reality provides inputs \\
Failure handling & Restart and try again & Must recover autonomously \\
Audience & Sympathetic experts & Indifferent users \\
Success metric & ``It works!'' & ``It keeps working.'' \\
\bottomrule
\end{tabular}
\caption{Structural differences between demo and deployment.}
\label{tab:demo-vs-deploy}
\end{table}

``Deploy or Die'' is not a rejection of the demo culture but its evolution. The demo proves that something is possible. Deployment proves that something is sustainable. The history of technology is littered with brilliant demos that never shipped --- and with mediocre demos that, through persistence of deployment, became infrastructure.

This thesis proposes that the transition from demo to deployment requires three pillars, each grounded in established research traditions.

% ==================================================================
% 3. PILLAR 1: THE SENSING-UNDERSTANDING-ADAPTING PIPELINE
% ==================================================================
\section{Pillar 1: The Sensing-Understanding-Adapting Pipeline}

\subsection{Affective Sensing as Foundation}

Picard's affective computing provides the first pillar's foundation. An autonomous system that cannot sense its environment is blind; one that can sense but not understand is overwhelmed; one that can understand but not adapt is brittle. The pipeline --- sense, understand, adapt --- is the minimal architecture for sustained autonomous operation.

Sensing in the affective computing tradition goes beyond physical measurement. It encompasses the recognition of emotional states, social signals, and contextual cues that determine how information should be interpreted. The same data --- a spike in user interaction frequency --- means different things in different affective contexts: excitement, anxiety, or compulsive behavior. Without affective sensing, the system cannot distinguish between these interpretations.

The sensor technologies for interactive environments tradition extends this to multi-sensor fusion: the integration of heterogeneous sensor streams into a unified perceptual model. The principle is that no single sensor is sufficient. Skin conductance alone cannot distinguish anger from excitement. Facial expression alone cannot distinguish genuine from performed emotion. Fusion --- the principled combination of multiple weak signals --- produces robust perception that no individual signal can achieve.

\subsection{Understanding Through Multiple Frames}

Understanding, in the Minsky tradition, is not a single process but a society of partial interpretations. \emph{The Society of Mind} \citep{minsky1986} proposes that intelligence arises from the interaction of many ``agents,'' each contributing a partial perspective. Applied to the sensing-understanding-adapting pipeline: a deployable system must interpret its sensor data through multiple frames simultaneously --- emotional, logical, contextual, historical --- and negotiate between them.

This is the architectural insight that separates demo systems from deployed systems. A demo can afford a single interpretation frame because its inputs are controlled. A deployed system faces ambiguous, contradictory, and adversarial inputs that require multiple simultaneous interpretations and a principled mechanism for selecting among them.

\subsection{Adaptation as Continuous Process}

Adaptation in the deployable system is not a one-time calibration but a continuous process. The system must adapt to:

\begin{itemize}[leftmargin=*, itemsep=2pt]
  \item \textbf{Distribution shift}: the real world changes, and yesterday's model is today's liability.
  \item \textbf{User evolution}: the people interacting with the system learn, change preferences, develop new needs.
  \item \textbf{Environmental degradation}: sensors drift, networks degrade, dependencies fail.
  \item \textbf{Self-modification}: the system's own outputs change its future inputs (feedback loops).
\end{itemize}

The adaptive systems tradition (reinforcement learning foundations) provides the tools: feedback loops, online learning, and the critical discipline of monitoring adaptation to ensure it converges rather than diverges.

% ==================================================================
% 4. PILLAR 2: BIO-INSPIRED RESILIENCE
% ==================================================================
\section{Pillar 2: Bio-Inspired Resilience}

\subsection{Biomimicry as Design Method}

The ``How to Grow (Almost) Anything'' tradition teaches that biological systems have solved many of the problems that engineered systems face --- but under constraints that engineers often ignore. Biological systems are:

\begin{itemize}[leftmargin=*, itemsep=2pt]
  \item \textbf{Self-repairing}: a cut heals, a broken bone knits, damaged neurons reroute.
  \item \textbf{Self-assembling}: embryonic development produces complex structures from simple rules.
  \item \textbf{Adaptive}: immune systems learn new threats without forgetting old ones.
  \item \textbf{Resource-efficient}: biological computation operates at orders of magnitude less energy than silicon computation.
\end{itemize}

These properties are not incidental --- they are the product of billions of years of evolutionary optimization under constraints remarkably similar to those facing deployed AI systems: limited resources, hostile environments, and the requirement for sustained operation without external maintenance.

\subsection{Sleep and Dream Architecture}

Perhaps the most striking bio-inspired principle for AI deployment comes from sleep neuroscience. The two-stage memory model \citep{marr1971, mcclelland1995} proposes that the brain maintains two learning systems: a fast learner (hippocampus) that quickly encodes new experiences, and a slow learner (neocortex) that gradually integrates them into long-term knowledge. The transfer between these systems occurs primarily during sleep, through mechanisms including sharp-wave ripples (compressed replay of recent experiences) and sleep spindles (thalamo-cortical oscillations that gate information into neocortical storage).

This architecture solves a fundamental problem in machine learning known as catastrophic forgetting: the tendency of neural networks to overwrite previously learned information when trained on new data. The biological solution --- Complementary Learning Systems (CLS) --- has been independently rediscovered in machine learning as experience replay \citep{lin1992, mnih2015} and elastic weight consolidation \citep{kirkpatrick2017}.

For deployable AI systems, the lesson is clear: continuous learning requires a sleep-like consolidation process. Systems that learn only during ``waking'' operation --- processing new data as it arrives --- will either forget what they knew or fail to learn from new experience. A dedicated consolidation phase, inspired by biological sleep, allows the system to integrate new learning without catastrophic forgetting, prune unnecessary connections, and stress-test its world model through free-running generation (the computational equivalent of dreaming).

\subsection{Self-Repair and Autonomous Recovery}

The ARISE principle (Autonomous Recovery through Intelligent Self-Evaluation) draws from both biological self-repair and space engineering's requirement for autonomous fault recovery. In orbital environments, there is no technician to call --- the system must diagnose and repair itself or fail permanently. This constraint, applied to AI deployment, produces systems that:

\begin{itemize}[leftmargin=*, itemsep=2pt]
  \item Monitor their own health metrics continuously.
  \item Detect degradation before it becomes failure.
  \item Execute recovery procedures without human intervention.
  \item Maintain a ``heartbeat'' that signals operational status.
\end{itemize}

The heart engine concept --- a self-regulating pacing mechanism that adjusts system behavior based on real-time health metrics --- is directly inspired by the biological cardiac system, which maintains stable operation through autonomous feedback without conscious control.

% ==================================================================
% 5. PILLAR 3: DEPLOYMENT AS DESIGN CONSTRAINT
% ==================================================================
\section{Pillar 3: Deployment as Design Constraint}

\subsection{The Deployment Paradox}

The communities that most need AI-driven solutions are often the communities with the least power to shape how those solutions are designed. This deployment paradox --- that the gap between designer and user is widest precisely where the stakes are highest --- is the central ethical challenge of AI deployment.

Space architecture offers a resolution: in orbital habitats, the crew IS the user. There is no gap between designer and deployer. The lesson for AI deployment: close the gap or close the project.

\subsection{Constraints as Liberation}

Space architecture teaches that extreme constraints do not limit design --- they redirect it. Zero gravity eliminates compression-based structures but enables tensegrity. Resource scarcity eliminates waste but enables closed-loop ecology. Radiation shielding requirements eliminate lightweight materials but enable multi-functional walls.

Applied to AI deployment:

\begin{itemize}[leftmargin=*, itemsep=2pt]
  \item Limited compute does not prevent intelligence --- it requires efficient intelligence.
  \item Limited data does not prevent learning --- it requires selective learning.
  \item Ethical constraints do not prevent capability --- they require responsible capability.
\end{itemize}

The ALARA principle (As Low As Reasonably Achievable) from radiation protection provides the general framework: perfection is not the standard; reasonable achievement given competing constraints is. This applies to fairness, accuracy, latency, and every other deployment metric. The engineer's job is not to eliminate tradeoffs but to make them explicit, documented, and defensible.

\subsection{Deployment Methodology}

The thesis proposes a deployment methodology based on the capstone experience:

\begin{enumerate}[leftmargin=*, itemsep=2pt]
  \item \textbf{Define deployment as a set of binary conditions.} ``Does it deploy?'' is answered by a checklist of specific, testable conditions --- not by subjective assessment.
  \item \textbf{Implement canary deployment.} Start in observation mode (log but do not act), then transition to active mode after pattern stabilization.
  \item \textbf{Design for graceful degradation.} When a component fails, the system should degrade to a simpler operating mode --- not crash.
  \item \textbf{Establish abort criteria before launch.} Know the conditions under which you stop, before you start.
  \item \textbf{Monitor continuously, iterate rapidly.} Post-deployment is not post-development.
\end{enumerate}

% ==================================================================
% 6. SERC FRAMEWORK
% ==================================================================
\section{SERC: A Social-Ethical-Regulatory-Cultural Evaluation Framework}

The three pillars --- sensing-understanding-adapting, bio-inspired resilience, and deployment as constraint --- address the technical dimensions of sustainable AI operation. But deployed AI systems operate in social contexts, and technical sustainability without social responsibility is insufficient.

The SERC framework evaluates AI systems across four dimensions:

\paragraph{Social.} Who is affected by this system? Who benefits? Who bears the costs? Are affected communities involved in design?

\paragraph{Ethical.} Does the system respect autonomy, prevent harm, ensure fairness, and maintain transparency? Which ethical framework (consequentialist, deontological, virtue, care) guides design decisions, and are the tradeoffs explicit?

\paragraph{Regulatory.} Does the system comply with relevant regulations (GDPR, EU AI Act, sector-specific requirements)? Is compliance achieved through design rather than post-hoc patching?

\paragraph{Cultural.} Does the system operate appropriately across cultural contexts? Does it encode cultural assumptions that may not transfer?

The capstone project --- SERC-Enhanced Adaptive Dreamwalk --- demonstrates that SERC evaluation can be integrated inline into an autonomous learning pipeline, rather than applied as a post-hoc audit. This is a significant departure from current practice, which treats ethical evaluation as an assessment performed after development rather than a constraint applied during development.

The inline approach has a specific advantage: it catches ethical issues at the point of learning, before they are consolidated into permanent system behavior. Post-hoc audits can identify problems but cannot prevent them from being learned in the first place.

% ==================================================================
% 7. CASE STUDY
% ==================================================================
\section{Case Study: SERC-Enhanced Adaptive Dreamwalk}

The capstone project integrates all three pillars:

\paragraph{Pillar 1 (Sensing-Understanding-Adapting).} The system senses incoming data through \texttt{feed\_collector}, understands it through dreamwalk's associative processing, and adapts through \texttt{brain\_compiler}'s weight updates. Multi-sensor fusion is applied through the combination of textual data, emotional weights (FIMP), and structural neuron relationships.

\paragraph{Pillar 2 (Bio-Inspired Resilience).} The dreamwalk pipeline is modeled on sleep architecture --- fast learning (data collection) followed by consolidation (dreamwalk replay and \texttt{brain\_compiler} integration). Selective consolidation based on emotional salience mirrors biological memory triage. The system includes autonomous health monitoring and graceful degradation.

\paragraph{Pillar 3 (Deployment as Constraint).} SERC evaluation is inline, not post-hoc. Deployment was defined by five binary conditions. Canary deployment (observation mode to active mode) was implemented. Abort criteria were established before launch.

Post-deployment metrics over 7 days:

\begin{table}[h]
\centering
\begin{tabular}{@{}ll@{}}
\toprule
\textbf{Metric} & \textbf{Result} \\
\midrule
Integration accuracy & 94.2\% (target: $>$90\%) \\
SERC flag rate & 3.1\% (normal range: 1--10\%) \\
System latency increase & +38\% (target: $<$50\%) \\
Uptime & 99.7\% \\
\bottomrule
\end{tabular}
\caption{Post-deployment metrics over 7 days.}
\label{tab:metrics}
\end{table}

The system met all deployment criteria and sustained operation without human intervention.

% ==================================================================
% 8. CONCLUSION
% ==================================================================
\section{Conclusion}

The evolution from ``Demo or Die'' to ``Deploy or Die'' reflects a maturation in the relationship between research and reality. Demonstrations prove possibility; deployments prove sustainability. The principles that enable this transition --- affective sensing, bio-inspired resilience, and deployment as design constraint --- are not post-hoc additions to the engineering process but foundational architectural decisions that must be present from conception.

Picard taught us that machines must sense affect to interact with humans. Minsky taught us that intelligence is a society of agents, not a monolithic process. Negroponte taught us that the interesting problems live at boundaries --- atoms and bits, technical and social, possible and sustainable. This thesis synthesizes these traditions into a framework for AI systems that ship: systems that not only work in the lab but survive, adapt, and operate responsibly in the world.

The SERC framework adds the dimension that the original research tradition underemphasized: social responsibility as an engineering constraint, not an afterthought. The inline approach --- evaluating ethics during learning rather than after deployment --- is, we argue, the natural extension of the ``deployment as constraint'' philosophy. If deployment is a design constraint, then ethical deployment is a design constraint. It belongs in the architecture, not in the documentation.

The final lesson is the simplest: Deploy or Die is not a threat. It is a design philosophy. It says: build things that work in the real world, that survive contact with reality, that adapt to conditions you did not anticipate, and that do so responsibly. It says: the world does not need more demos. It needs more deployments.

% ==================================================================
% REFERENCES
% ==================================================================
\bibliographystyle{plainnat}

\begin{thebibliography}{20}

\bibitem[Adamson and Smith(2018)]{adamson2018}
Adamson, A.~S. and Smith, A. (2018).
\newblock Machine learning and health care disparities in dermatology.
\newblock \emph{JAMA Dermatology}, 154(11):1247--1248.

\bibitem[Chouldechova(2017)]{chouldechova2017}
Chouldechova, A. (2017).
\newblock Fair prediction with disparate impact: A study of bias in recidivism prediction instruments.
\newblock \emph{Big Data}, 5(2):153--163.

\bibitem[Crick and Mitchison(1983)]{crick1983}
Crick, F. and Mitchison, G. (1983).
\newblock The function of dream sleep.
\newblock \emph{Nature}, 304(5922):111--114.

\bibitem[Hobson and Friston(2012)]{hobson2012}
Hobson, J.~A. and Friston, K.~J. (2012).
\newblock Waking and dreaming consciousness: Neurobiological and functional considerations.
\newblock \emph{Progress in Neurobiology}, 98(1):82--98.

\bibitem[Hopfield(1982)]{hopfield1982}
Hopfield, J.~J. (1982).
\newblock Neural networks and physical systems with emergent collective computational abilities.
\newblock \emph{Proceedings of the National Academy of Sciences}, 79(8):2554--2558.

\bibitem[Kirkpatrick et~al.(2017)]{kirkpatrick2017}
Kirkpatrick, J., et~al. (2017).
\newblock Overcoming catastrophic forgetting in neural networks.
\newblock \emph{Proceedings of the National Academy of Sciences}, 114(13):3521--3526.

\bibitem[Kleinberg et~al.(2016)]{kleinberg2016}
Kleinberg, J., Mullainathan, S., and Raghavan, M. (2016).
\newblock Inherent trade-offs in the fair determination of risk scores.
\newblock \emph{arXiv:1609.05807}.

\bibitem[Lin(1992)]{lin1992}
Lin, L.~J. (1992).
\newblock Self-improving reactive agents based on reinforcement learning, planning and teaching.
\newblock \emph{Machine Learning}, 8(3):293--321.

\bibitem[Marr(1971)]{marr1971}
Marr, D. (1971).
\newblock Simple memory: A theory for archicortex.
\newblock \emph{Philosophical Transactions of the Royal Society B}, 262(841):23--81.

\bibitem[McClelland et~al.(1995)]{mcclelland1995}
McClelland, J.~L., McNaughton, B.~L., and O'Reilly, R.~C. (1995).
\newblock Why there are complementary learning systems in the hippocampus and neocortex.
\newblock \emph{Psychological Review}, 102(3):419--457.

\bibitem[Minsky(1986)]{minsky1986}
Minsky, M. (1986).
\newblock \emph{The Society of Mind}.
\newblock Simon \& Schuster.

\bibitem[Mnih et~al.(2015)]{mnih2015}
Mnih, V., et~al. (2015).
\newblock Human-level control through deep reinforcement learning.
\newblock \emph{Nature}, 518(7540):529--533.

\bibitem[Negroponte(1995)]{negroponte1995}
Negroponte, N. (1995).
\newblock \emph{Being Digital}.
\newblock Alfred A.~Knopf.

\bibitem[Picard(1995)]{picard1995}
Picard, R.~W. (1995).
\newblock Affective computing.
\newblock \emph{Technical Report No.~321, Perceptual Computing Section}.

\bibitem[Picard(1997)]{picard1997}
Picard, R.~W. (1997).
\newblock \emph{Affective Computing}.
\newblock MIT Press.

\bibitem[Voss et~al.(2009)]{voss2009}
Voss, U., et~al. (2009).
\newblock Lucid dreaming: A state of consciousness with features of both waking and non-lucid dreaming.
\newblock \emph{Sleep}, 32(9):1191--1200.

\bibitem[Winner(1980)]{winner1980}
Winner, L. (1980).
\newblock Do artifacts have politics?
\newblock \emph{Daedalus}, 109(1):121--136.

\bibitem[Zuboff(2019)]{zuboff2019}
Zuboff, S. (2019).
\newblock \emph{The Age of Surveillance Capitalism}.
\newblock PublicAffairs.

\end{thebibliography}

\end{document}
